% An auxiliary method obstruction, not a resolution of the C.4 residual case.
% Owned by /root/frobenius; C.4y is a proved method limitation, not main-gate movement.
\paragraph{C.4y (Proved): a parity obstruction for semihomogeneous Chern constructions.}
Let $A/\mathbb C$ be an abelian surface with
$\operatorname{End}(A)=\mathbb Z$ and a primitive polarization $L$ of type
$(1,2)$. Put $X=A\times\widehat A$, and write
\[
 \ell_1=p_A^*c_1(L),\qquad
 \ell_2=p_{\widehat A}^*c_1(\widehat L),\qquad
 P=c_1(\mathcal P_A),
\]
where $\widehat L$ is the dual polarization, so that
$\lambda_{\widehat L}\lambda_L=[2]$. Let $R_{\rm sh}$ be the integral
graded subring of numerical cycle classes on $X$ generated by the Chern
classes of all semihomogeneous vector bundles on $X$. Then
\[
                    P^2/2\notin R_{\rm sh}^{2}.
\]
Moreover, the mixed Matsusaka coefficient of every class in
$R_{\rm sh}^{3}$ is even. Here we normalize the mixed coefficient to
take the value $1$ on $P^3/6$; an overall change of Poincar\'e sign does not
change the parity assertion. These statements rule out this construction
of the missing class in C.4r. They assert nothing about arbitrary algebraic
cycles or arbitrary vector bundles on $X$.

\emph{Proof.}
The hypothesis on endomorphisms gives
$\operatorname{NS}(X)=\mathbb Z\ell_1\oplus\mathbb ZP\oplus
\mathbb Z\ell_2$, as in C.4t. In an integral symplectic basis we may use
the alternating forms
\[
 \ell_1=a_1a_2+2a_3a_4,\qquad
 \ell_2=2b_1b_2+b_3b_4,\qquad
 P=\sum_{i=1}^4a_ib_i,
\]
where juxtaposition means exterior product. Direct expansion gives
\begin{equation}\label{eq:c4-sh-forms}
 \chi(a\ell_1+bP+c\ell_2)=(2ac-b^2)^2,
 \qquad P^3=-3\ell_1\ell_2P.
\end{equation}
The six degree-two monomials in $\ell_1,P,\ell_2$ are numerically
independent: their intersection matrix, in the order
$(\ell_1^2,\ell_1P,\ell_1\ell_2,P^2,P\ell_2,\ell_2^2)$, is
\[
 \begin{pmatrix}
 0&0&0&0&0&16\\
 0&0&0&0&-8&0\\
 0&0&16&-8&0&0\\
 0&0&-8&24&0&0\\
 0&-8&0&0&0&0\\
 16&0&0&0&0&0
 \end{pmatrix}.
\]
In particular, extraction of the $P^2$ coefficient is well defined on
their rational span; we do not need to identify numerical and
homological equivalence for all cycles on $X$.

Let $E$ be semihomogeneous of rank $r>0$ and let
$D=c_1(E)=a\ell_1+bP+c\ell_2$. We have
\begin{equation}\label{eq:c4-sh-chern}
 \operatorname{ch}(E)=r\exp(D/r),\quad
 c_2(E)=\frac{r-1}{2r}D^2,\quad
 c_3(E)=\frac{(r-1)(r-2)}{6r^2}D^3.
\end{equation}
For completeness, the full semihomogeneous statement follows from
\citet[Theorems 2.4 and 2.6, p.~6, and Lemma 2.9, pp.~6--7]
{frGrossKaurUlirschWerner2026}. A simple semihomogeneous bundle is
$f_*M$ for an isogeny $f$ and a line bundle $M$; after pullback it is the
sum of the translates of $M$, which have the same numerical Chern class.
The injectivity of rational pullback under a finite isogeny proves the
exponential formula. Their Theorem 2.4 writes a general such bundle as a
sum of simple bundles of the same slope, tensored by unipotent bundles.
A unipotent bundle has a filtration with trivial line quotients, and
therefore changes the Chern character only by its rank. This proves
\eqref{eq:c4-sh-chern} without assuming that $E$ is simple.

Restrictions of $c_2(E)$ to $A\times\{0\}$ and
$\{0\}\times\widehat A$, together with Hirzebruch--Riemann--Roch on
the abelian fourfold, give the three integral quantities
\begin{equation}\label{eq:c4-sh-integrality}
 \frac{2(r-1)a^2}{r},\qquad
 \frac{2(r-1)c^2}{r},\qquad
 \chi(E)=\frac{(2ac-b^2)^2}{r^3}.
\end{equation}
All assertions below concern the prime $2$ only. Write
$\mathbb Z_{(2)}=\{u/v\in\mathbb Q:2\nmid v\}$.
If $r$ is odd, the coefficient of $P^2$ in $c_2(E)$, and the
coefficients of $\ell_1P$, $P\ell_2$, and $\ell_1\ell_2$, already
belong to $\mathbb Z_{(2)}$.

Suppose $s=v_2(r)>0$, and put $u=v_2(a)$, $v=v_2(b)$ and
$w=v_2(c)$, allowing the value $+\infty$ for a zero coefficient.
The first two expressions in \eqref{eq:c4-sh-integrality} imply
\[
                          2u+1\ge s,\qquad 2w+1\ge s.
\]
If $2v\le s$, then $v_2(2ac)>2v$: for even $s$, we have
$u,w\ge s/2$; for odd $s$, we have
$1+u+w\ge s>2v$. Consequently
$v_2(2ac-b^2)=2v$, and the last expression in
\eqref{eq:c4-sh-integrality} has valuation
$4v-3s\le-s<0$, a contradiction. We have proved
\begin{equation}\label{eq:c4-sh-b}
                                2v>s.
\end{equation}
Thus
\[
             \frac{(r-1)b^2}{2r}\in\mathbb Z_{(2)}.
\]
This is exactly the $P^2$ coefficient of $c_2(E)$.
The two mixed coefficients $(r-1)ab/r$ and $(r-1)bc/r$ are also
$2$-integral, since $u+v,w+v\ge s$. Finally $u+w\ge s$:
the restriction inequalities allow a failure only when $s$ is odd and
$u=w=(s-1)/2$; but then \eqref{eq:c4-sh-b} gives
$v_2(2ac-b^2)=s$, contradicting $2s-3s\ge0$ in
\eqref{eq:c4-sh-integrality}. Hence $(r-1)ac/r$ is $2$-integral.

Every degree-two element of $R_{\rm sh}$ is an integral sum of second
Chern classes and products of first Chern classes. The preceding
calculation shows that its $P^2$ coefficient lies in
$\mathbb Z_{(2)}$, whereas the coefficient for $P^2/2$ is $1/2$.
Numerical independence of the six monomials proves the first assertion.

For the degree-three assertion, the only bidegree $(3,3)$ term of
$D^3$ is
\[
 b^3P^3+6abc\ell_1\ell_2P
       =-b(2ac-b^2)P^3.
\]
Thus the mixed coefficient of $c_3(E)$, normalized by $P^3/6$, is
\begin{equation}\label{eq:c4-sh-third}
                  -\frac{(r-1)(r-2)b(2ac-b^2)}{r^2}.
\end{equation}
For odd $r$ this is in $2\mathbb Z_{(2)}$. For even $r$, set
$h=v_2(2ac-b^2)$; we may disregard the zero expression. The Euler
characteristic in \eqref{eq:c4-sh-integrality} gives
$h\ge\lceil3s/2\rceil$, and \eqref{eq:c4-sh-b} gives
$v\ge\lfloor s/2\rfloor+1$. If $s\ge2$, the valuation of
\eqref{eq:c4-sh-third} is at least
\[
 1+\lfloor s/2\rfloor+1+\lceil3s/2\rceil-2s=2.
\]
If $s=1$, then $r=2$ makes $c_3(E)=0$; otherwise
$v_2(r-2)\ge2$, and the valuation is again positive.

For a product of $c_2(E)$ with any integral divisor, its relevant
bidegree $(3,3)$ coefficient is an integral linear combination of
$6\kappa$ and $-2\tau$, where $\kappa$ is the $P^2$ coefficient and
$\tau$ is one of the other three mixed coefficients just proved to
lie in $\mathbb Z_{(2)}$. Its mixed coefficient is therefore in
$2\mathbb Z_{(2)}$. The same conclusion for a product of three
divisors follows directly from \eqref{eq:c4-sh-forms}. These three
types generate $R_{\rm sh}^{3}$, so every mixed coefficient is
$2$-integral and even. Finally the mixed coefficient of any actual
integral curve is an integer: it is the scalar of its mixed integral
homomorphism, under the integral correspondence construction of C.4i,
and $\operatorname{End}(A)=\mathbb Z$. Hence it lies in $2\mathbb Z$.
\hfill$\square$

\emph{Scope.}
The argument also covers Chern polynomials of virtual classes generated
additively by semihomogeneous bundles, since inverse total Chern classes
are integral polynomials in their components. More explicitly, let
$q:A_y\times A_a\times\widehat A_\xi\to A_a\times\widehat A_\xi$.
The diagonal ideal exact sequence gives
\[
 [Rq_*(I_{\Delta_{y,a}}\otimes L_y\otimes\mathcal P_{y,\xi})]
  =[p_{\widehat A}^*\Phi_{\mathcal P}(L)]
       -[p_A^*L\otimes\mathcal P_A]
\]
in integral $K_0(X)$. The Fourier transform of $L$ is a vector bundle,
and its pullback under $\lambda_L$ is
$H^0(A,L)\otimes L^{-1}$, by the theorem-of-the-square identity for
the pulled-back Poincar\'e bundle and cohomology of an ample line
bundle. It is consequently semihomogeneous (alternatively the same
pullback already proves the needed exponential Chern identity).
Thus this familiar odd-rank virtual construction lies in the excluded
ring as well. We have not proved that arbitrary Fourier--Mukai or
moduli kernels belong to that ring. In particular, the argument does
not decide the numerical integrality of $P^2/2$ for $X$.
